\documentclass[10pt,a4paper]{article}

% ---------- packages ----------
\usepackage[top=1.8cm,bottom=1.8cm,left=1.8cm,right=1.8cm]{geometry}
\usepackage{multicol}
\usepackage{booktabs}
\usepackage{tabularx}
\usepackage{xcolor}
\usepackage{titlesec}
\usepackage{enumitem}
\usepackage{microtype}
\usepackage{parskip}
\usepackage{lmodern}
\usepackage[T1]{fontenc}
\usepackage[utf8]{inputenc}
\usepackage{hyperref}
\usepackage{graphicx}

% ---------- colours ----------
\definecolor{tuwblue}{RGB}{0,102,153}
\definecolor{lightgray}{RGB}{245,245,245}
\definecolor{darkgray}{RGB}{80,80,80}

% ---------- section formatting ----------
\titleformat{\section}{\normalsize\bfseries\color{tuwblue}}{}{0em}{}[\vspace{-4pt}\hrule\vspace{4pt}]
\titlespacing*{\section}{0pt}{6pt}{2pt}

% ---------- list spacing ----------
\setlist[itemize]{noitemsep, topsep=1pt, leftmargin=*}

\pagestyle{empty}

% ============================================================
\begin{document}

% ---------- header ----------
\begin{center}
  {\large\bfseries\color{tuwblue}
    Can Wrist Motion Reveal How You Carry a Box?}\\[3pt]
  {\small\color{darkgray}
   Management Summary \quad|\quad TU Wien ASM2 \quad|\quad
   Viktor Hoffmann \quad|\quad May 2026}
\end{center}
\vspace{2pt}
\hrule height 1.2pt
\vspace{6pt}

% ---------- two-column body ----------
\begin{multicols}{2}

% ~~~~~~~~~~~~~~~~~~~~~~~~~~~~~~~~~~~~~~~~~~
\section{The Question}

When a person carries a box with one hand, their two wrists move differently:
the active wrist is constrained by the load while the passive wrist swings freely.
We asked whether this left-right asymmetry is strong enough that an
\textbf{unsupervised clustering algorithm} --- given no labels --- would spontaneously
group windows of sensor data by carrying condition.
Two wrist-mounted IMUs (left and right) recorded three-axis gyroscope and
accelerometer signals from two subjects across 2 box sizes, 2 surfaces, and
3 carrying conditions (left hand, right hand, both hands) at $\approx$100\,Hz,
yielding 47{,}413 raw samples.

% ~~~~~~~~~~~~~~~~~~~~~~~~~~~~~~~~~~~~~~~~~~
\section{Step 1: Start Simple --- One Wrist at a Time}

We began with a \textbf{single-sensor} analysis: extract 91 time-domain features
per 1-second window from one wrist (mean, std, RMS, kurtosis, jerk magnitude,
etc.\ across 10 signal channels), scale with RobustScaler, and cluster.
By construction this approach cannot represent cross-wrist asymmetry.

The clusters were clear and interpretable, but they captured the \textbf{wrong
thing}: the dominant structure was \textit{motion intensity}.
High-kurtosis, high-jerk windows (picking up, setting down) separated from
low-intensity walking windows.
All carrying-condition ARIs were $\leq$\,0.08.
\textit{Wrist motion intensity varies by movement phase, not by which hand
holds the box.}

% ~~~~~~~~~~~~~~~~~~~~~~~~~~~~~~~~~~~~~~~~~~
\section{Step 2: Add the Second Wrist}

Single-sensor analysis cannot test the lateralisation hypothesis.
We built a \textbf{paired} feature matrix: one row combines both wrists
simultaneously, adding 188 cross-wrist asymmetry features
(absolute difference $|f_R - f_L|$, ratio $f_R / f_L$, and Pearson correlation
per axis), giving \textbf{370 features} per window.

Initial paired results looked suspiciously strong: $k$\,=\,2 silhouette\,=\,0.78
pooled --- far above single-sensor.
Yet all carrying-condition ARIs remained $\approx$\,0.
Inspecting the scatter revealed the truth: a handful of extreme windows formed
one cluster on their own, far removed from all other data.

% ~~~~~~~~~~~~~~~~~~~~~~~~~~~~~~~~~~~~~~~~~~
\section{Step 3: The Artefact, and How We Fixed It}

Examining which features separated the outlier cluster, we found extreme
raw-sensor kurtosis values (up to 97; dataset 99th percentile: 15).
These are \textbf{impact windows} --- moments when one wrist made accidental
contact with a surface.
Because they generate extreme asymmetry-feature values, they dominated the
entire paired clustering.

We removed windows where any raw L/R sensor kurtosis exceeded 30, while
deliberately \textit{excluding} the new jerk-channel kurtosis features
(Ajerk, Gjerk are impulsive by construction; their kurtosis is naturally
high during normal motion and would falsely flag the majority of windows).
This targeted filter removed exactly \textbf{14 windows (2.3\,\%)} while
preserving 97.7\,\% of the data.

% [INSERT FIGURE: data/features/plots/Viktor/pass\_b/scatter\_one\_handed\_carry\_pass\_b.png
%  Caption: Ground-truth one\_handed\_carry label overlay in classifier feature space
%  (Viktor, Pass B). A spatial gradient is visible: both-handed windows concentrate
%  upper-left, one-handed windows spread lower-right --- motivating the clustering.]

% ~~~~~~~~~~~~~~~~~~~~~~~~~~~~~~~~~~~~~~~~~~
\section{Step 4: Rethinking the Pipeline --- Removing PCA}

With clean data, the previous pipeline (RobustScaler\,$\to$\,PCA\,$\to$\,cluster)
still failed: silhouette remained $\approx$\,0.55--0.83 but external ARI stayed near
zero.
Two structural problems were identified:

\begin{itemize}
  \item \textbf{PCA maximises variance, not class separability.}  The dominant
        variance in 370 features comes from motion intensity (the same
        phase-correlated structure found in single-sensor analysis), not from
        the carrying-condition signal.  PCA projects the discriminative direction
        into low-variance components and discards it.
  \item \textbf{Feature redesign.}  The old asymmetry set included
        \texttt{symm\_idx} (numerically unstable near zero) and
        \texttt{log\_ratio} (loses sign).  These were replaced by a linear
        \texttt{ratio}\,$= f_R / (f_L + \varepsilon)$, and a dedicated gyro jerk
        channel (Gjerk) was added alongside accel jerk (Ajerk).
\end{itemize}

The new pipeline \textbf{removes PCA from clustering entirely}.
RobustScaler + $\pm$50 clip is applied, then algorithms run in the full
scaled feature space.
A 2-component PCA is computed \textit{only} for scatter plot visualisation
(not used for clustering).

% ~~~~~~~~~~~~~~~~~~~~~~~~~~~~~~~~~~~~~~~~~~
\section{Step 5: Two-Pass Evaluation}

We ran two passes per subject, each stratified within-subject to eliminate
between-subject confounding:

\textbf{Pass A --- full 370-feature space.}
K-Means, Ward, and GMM (diagonal covariance, $n$\,=\,2--5).
DBSCAN skipped: $\varepsilon$ is not interpretable in $\sim$370 dimensions.
Result: silhouette\,=\,0.81--0.83 (very high internally consistent structure),
but \textit{all} carrying-condition ARIs\,$\approx$\,0.
The full feature space clusters strongly --- but by motion dynamics and
session-level variation, not by carrying condition.

\textbf{Pass B --- classifier top-10 features.}
A colleague's supervised Random Forest (F1\,$\approx$\,0.88--0.92) identified the
10 most discriminative features: mean accelerations per wrist
(\texttt{L/R\_AX/AY/AZ\_mean}) and four cross-wrist asymmetry features.
We clustered \textit{only} in that 10-dimensional space.
DBSCAN is included here (10 dims: $\varepsilon$ is interpretable).
Results:

\begin{itemize}
  \item \textbf{Viktor, $k$\,=\,2:}
        ARI(\texttt{one\_handed\_carry})\,=\,\textbf{0.40}.
        Cluster\,0 (129 windows): \textbf{92.2\,\% both-handed}.
        Cluster\,1 (196 windows): \textbf{75.0\,\% one-handed}.
        The binary carrying split is clearly recovered without labels.
  \item \textbf{Viktor, GMM $n$\,=\,5:}
        ARI(\texttt{carrying\_hand})\,=\,\textbf{0.52}.
        Three of five clusters are perfectly pure:
        107 windows at \textbf{100\,\% both-handed},
        58 windows at \textbf{100\,\% right-handed},
        45 windows at \textbf{100\,\% left-handed}.
        The three carrying conditions emerge as natural density modes.
  \item \textbf{David, $k$\,=\,3:}
        ARI(\texttt{one\_handed\_carry})\,=\,\textbf{0.35}.
        Weaker but real; David's signal requires $k \geq 3$ because his
        wrist kinematics differ less sharply between left- and right-hand trials.
\end{itemize}

DBSCAN collapsed to 1 cluster for both subjects (with 8--16\,\% noise), confirming
no compact density-gap between conditions --- only distributional differences that
K-Means and GMM exploit through centroid geometry.

% [INSERT FIGURE: data/features/plots/Viktor/pass\_b/scatter\_cluster\_kmeans\_k2\_pass\_b.png
%  Caption: Viktor, Pass B, K-Means $k$\,=\,2 in 2D PCA view of the 10-feature space.
%  Blue (cluster 0) concentrates upper-left = both-handed; orange (cluster 1) spreads
%  right and lower = one-handed. ARI = 0.40.]

% [INSERT FIGURE: data/features/plots/Viktor/pass\_b/scatter\_cluster\_gmm\_n5\_pass\_b.png
%  Caption: Viktor, Pass B, GMM $n$\,=\,5. Blue = 100\,\% both-handed, orange =
%  100\,\% right-handed, purple = 100\,\% left-handed. Three carrying conditions
%  emerge as natural Gaussian density modes. CH-ARI = 0.52.]

% [INSERT FIGURE: data/features/plots/David/pass\_b/scatter\_cluster\_kmeans\_k2\_pass\_b.png
%  Caption: David, Pass B, K-Means $k$\,=\,2. The geometric split (left vs.\ right)
%  is visible, but both clusters are majority one-handed --- David requires $k \geq 3$
%  for condition-aligned structure (OHC-ARI = 0.35 at $k$\,=\,3).]

% ~~~~~~~~~~~~~~~~~~~~~~~~~~~~~~~~~~~~~~~~~~
\section{Conclusions}

\renewcommand{\arraystretch}{1.1}
\begin{tabularx}{\columnwidth}{@{}lX@{}}
\toprule
\textbf{Hypothesis} & \textbf{Verdict} \\
\midrule
H1 --- one-sided vs.\ two-handed separable &
  \textbf{Supported (Viktor).}
  K-Means $k$\,=\,2 ARI\,=\,0.40; cluster purities 92\,\% / 75\,\%.
  Weaker for David (ARI\,=\,0.35 at $k$\,=\,3). \\[3pt]
H2 --- carrying hand (L/R/both) separable &
  \textbf{Strongly supported (Viktor).}
  GMM $n$\,=\,5 recovers all three conditions with CH-ARI\,=\,0.52;
  three clusters at 100\,\% purity. \\[3pt]
H3 --- motion phases distinguishable &
  Weakly supported (max phase ARI\,$\leq$\,0.16).
  Phase correlates with intensity, not carrying condition. \\
\bottomrule
\end{tabularx}

\vspace{4pt}

The central lessons are methodological.
\textbf{First, PCA is the wrong dimensionality reduction for this task}: it
finds directions of maximum variance (motion intensity), not class separability,
discarding the carrying-condition signal in the process.
\textbf{Second, feature selection matters more than feature count}: 10 features
from a supervised classifier outperform 370 features fed through PCA by a wide
margin.
\textbf{Third, filter design requires domain knowledge}: the kurtosis filter must
exclude jerk channels (impulsive by construction) or it over-removes 55\,\% of
the data, fatally weakening every subsequent result.

Two actionable conclusions emerge: (1) mean wrist accelerations and four
cross-wrist asymmetry features are sufficient to cluster carrying conditions
without labels; (2) a supervised feature-selection step followed by unsupervised
clustering is a viable pipeline for future studies with $n \geq 10$ subjects,
where generalisation can be properly evaluated.

\end{multicols}

\vspace{4pt}
\hrule
\vspace{3pt}
{\footnotesize\color{darkgray}
Plots: \texttt{data/features/plots/\{subject\}/\{pass\}/}.\quad
Full metrics: \texttt{data/features/metrics\_summary.csv}.\quad
Documentation: \texttt{clustering.md}.}

\end{document}
